\title{RingRail: A Self-Identifying, Self-Synchronizing Ring-Bus Communication Protocol for FPGAs in Medium-Voltage Applications}
\subtitle{GridTran · Atlanta, GA}
\period{2026}
\date{2026-04}
\order{2}
\tags{control, patent}

\attach{Slides}{/assets/ringrail_slides.pdf}

\begin{document}

RingRail is the fibre control bus for a modular medium-voltage converter. One
main FPGA commands a series stack of bridge boards, each with its own FPGA
floating at a different potential, over optical rings. A board learns its
address from where it sits in the ring, and every board switches on the same
edge however far down the ring it is.

\section{How it works}

\begin{itemize}
  \item \textbf{Two trains per ring.} A slow train carries CRC-checked command
        frames out and status back. A fast train carries only a timing tick.
        The slow train says what to do, the fast train says when.
  \item \textbf{Self-identifying.} A board's index is its position in the ring,
        picked up from the traffic. There is no address switch, every board of
        a class runs the same image, and two boards can be swapped without
        reconfiguring anything.
  \item \textbf{Self-synchronizing.} Each board delays its trigger by the hops
        still ahead of the pulse, so all of them fire together, and a local
        phase-locked loop tracks the main FPGA across crystal offsets and
        rides over a lost tick.
  \item \textbf{Fail-safe.} The main FPGA watches its own frames come back
        around the ring. A broken ring or a lost timebase takes every board to
        a safe state.
\end{itemize}

\section{Patent}

D.~Divan, J.~Benzaquen, P.~R.~Nallamatti, S.~Manjrekar and S.~Itum,
``Self-Identifying, Self-Synchronizing Ring-Bus Communication Protocol for
FPGAs in Medium-Voltage Applications,'' provisional patent disclosure,
GridTran, April 2026.

\end{document}
